\section{Across horizons, kinds of market and seeds}
\label{ch:horizons}

The previous sections built and assessed the families mostly at one hour.  This section
reads the same families across the five horizons of the study, one minute, fifteen minutes,
one hour, six and twelve hours, because the horizon changes what the task is and which family
can do it; then by kind of market; then for robustness, across seeds, calibration and run
length.  The locked test set and the return to the four questions follow in
\cref{sec:horizons:test,sec:horizons:questions}.

\subsection{Short horizons: one and fifteen minutes}
\label{sec:horizons:short}
\label{sec:exp:short}


At a horizon of one minute the task changes: 89\% of the labels are \textsc{const}, the
direction of a move is nearly unpredictable and the whole gain is in whether the mid will
move (\cref{sec:expl:classes,sec:expl:activity}).  The early LOB-only models with undersampling scored $0.550$
at 60~s against $0.536$ for Mean Reversion, $0.532$ at 120~s and $0.518$ at 300~s; the
numbers below are from the final configuration of \cref{sec:stateful:hand} (five
lookbacks, hierarchical memory, time-first encoder, hand features, two passes), trained at
60~s and 900~s.
The controls of \cref{tab:exp:nostate:horizons} sample one tensor row in twenty at the
row's time, with the label of the last tensor row before $t+H$; the deep runs sample every
market of every block at the block's time, with the label of the mid-price index.  At one
hour the two populations coincide, because the control was built from the run's own
predictions; at one minute only 27 of the 598{,}872 tensor rows coincide with a run sample.
The control was therefore rebuilt on the runs' own validation samples at each horizon (9.9~M samples at both
60 and 900~s), with the same features, learner and protocol as in \cref{sec:stateless:control}.
Every model then receives the same two cross-fitted logit offsets of
\cref{sec:models:protocol}, with the grid widened to $\pm 2$: at one minute both learners
sit at the edge of the $\pm 0.6$ grid used at one hour, and the offsets are worth $+5.4$
points to LightGBM ($0.6086 \to 0.6623$) and $+5.8$ to the network.  Any one-minute number
without this calibration compares decision thresholds, not models; at fifteen minutes the
same offsets are worth $+0.5$, and at one hour $+0.3$.

\begin{table}[t]
\centering
\caption{Short horizons against one hour: the no-state control, the tabular network on the
control's table and the stateful model, every one with its own two cross-fitted logit
offsets (DOWN, UP), on the runs' validation samples of 13 March.  Balanced accuracy after
calibration; $\Delta$ in points against the calibrated LightGBM with 95\% hour-block
bootstrap intervals.  Labels DOWN / CONST / UP: $5.7 / 89.1 / 5.2\%$ at 1~min,
$16.4 / 69.0 / 14.6\%$ at 15~min.}
\label{tab:exp:short}
\small
\begin{fittable}\begin{tabular}{llcccc}
\toprule
Horizon & Model & All & $\Delta$ [95\% CI] & Dedup. & $\Delta$ [95\% CI] \\
\midrule
1 min  & LightGBM G5 (offsets 1.6, 1.6)  & 0.6623 & --- & 0.6143 & --- \\
       & TabM, 3 seeds                   & 0.6656 & $+0.33$ [$+0.26$, $+0.39$] & 0.6133 & $-0.10$ [$-0.31$, $+0.10$] \\
       & Stateful, 2 passes              & 0.6685 & $+0.62$ [$+0.46$, $+0.82$] & 0.6317 & $+1.74$ [$+1.31$, $+2.22$] \\
       & Stateful, 2 passes, residual    & \textbf{0.6760} & $\mathbf{+1.37}$ [$+1.16$, $+1.58$] & \textbf{0.6346} & $\mathbf{+2.04}$ [$+1.76$, $+2.30$] \\
\midrule
15 min & LightGBM G5 (offsets 0.4, 0.4)  & 0.6213 & --- & 0.6197 & --- \\
       & TabM, 3 seeds                   & \textbf{0.6235} & $+0.22$ [$+0.07$, $+0.36$] & \textbf{0.6248} & $+0.51$ [$+0.35$, $+0.68$] \\
       & Stateful, 2 passes              & 0.6206 & $-0.07$ [$-0.34$, $+0.19$] & 0.6158 & $-0.39$ [$-0.63$, $-0.13$] \\
\midrule
1 h    & LightGBM G5 (offsets 0.2, 0.3)  & 0.6044 & --- & 0.5943 & --- \\
       & TabM, 3 seeds                   & 0.6033 & $-0.11$ [$-0.29$, $+0.06$] & 0.5922 & $-0.21$ [$-0.45$, $+0.00$] \\
       & Stateful, 2 passes              & 0.5902 & $-1.42$ [$-1.83$, $-1.00$] & 0.5797 & $-1.46$ [$-1.76$, $-1.14$] \\
       & Stateful, 2 passes, residual    & \textbf{0.6048} & $+0.04$ [$-0.18$, $+0.24$] & 0.5933 & $-0.10$ [$-0.26$, $+0.05$] \\
\bottomrule
\end{tabular}\end{fittable}
\end{table}

\Cref{tab:exp:short} gives the result, and it reverses with the horizon.  At one minute the
stateful model without trees is $+0.62$ [$+0.46$, $+0.82$] above the calibrated trees on
all samples and $+1.74$ [$+1.31$, $+2.22$] on the deduplicated ones, ahead on every subset
($+0.45$ at extreme mids, $+0.64$ central).  TabM on the same sixty-one features and the
same samples, with the same calibration, is $+0.33$ on all samples and $-0.10$ on the
deduplicated ones: half of the gain on all samples is the learner, and the whole of the
gain on the sparse markets is the state, $+1.84$ over TabM.  Given the cross-fitted trees as
a base, the residual model is the best model of this thesis at one minute, $+1.37$
[$+1.16$, $+1.58$] on all samples and $+2.04$ [$+1.76$, $+2.30$] on the deduplicated ones,
above the calibrated trees on every subset and $+1.0$ to $+2.1$ above TabM.  At fifteen
minutes the stateful model ties the calibrated trees on all samples ($-0.07$ [$-0.34$,
$+0.19$]) and is below them on the deduplicated ($-0.39$) and extreme ($-0.27$) ones, while
TabM is $+0.22$ and $+0.51$: against the same learner the online, time-ordered regime costs
$-0.3$ on all samples and $-0.9$ on the sparse ones.  At one hour it costs $-1.4$ and the
residual model only recovers it (\cref{sec:stateful:calibration}).  Before calibration the
ordering is the same but the gaps are not: the raw stateful model is $+0.23$ (not
significant) at one minute and $-0.76$ [$-1.08$, $-0.45$] at fifteen, and the prior
correction of \cref{sec:stateful:hand} is worse than nothing at fifteen minutes
($0.6084 \to 0.5587$).

The reading is that the window of book events, the five lookbacks and the memory carry
information about the next minute of a market that the sixty-one summary features throw
away, chiefly on the quiet markets where a summary of the last hour says little about the
next sixty seconds, and that by fifteen minutes the summary is complete and the sequential
regime only costs.  The gain is small in absolute terms, under two points on a scale where
the activity clock alone is worth seven, but it is the one place in this study where the
stateful model beats a calibrated tree model on the same inputs with an interval that
excludes zero on every subset.  The residual model was not run at fifteen minutes.

\subsection{One hour}
\label{sec:horizons:1h}

One hour is the horizon at which most of the design was done, and its result is the tie of
\cref{sec:stateful:calibration}: calibrated against calibrated, the trees on sixty-one
hand-made features (0.6044 on all samples, 0.5943 deduplicated), TabM on the same table
($-0.11$), and the two-pass residual stateful model ($+0.04$) are within their intervals of
one another, the stateful model without trees is $-1.4$, and the best LOB-only model without
the hand features is four points below.  Half of the labels are \textsc{const}; the activity
clock decides whether the price moves and the reversal of the recent change decides which
way; the combination of the stateful model with the trees adds two to three tenths of a
point, most of it on the sparse markets.  The one-hour task is a summary task: everything a
model needs is in a few dozen numbers computed at the sample time.

\subsection{Six and twelve hours}
\label{sec:horizons:long}

At six and twelve hours the \textsc{const} class shrinks to a third of the samples and the
direction is the task.  The LOB-only stateful model of \cref{sec:stateful:chain} is
$+1.5$ above Mean Reversion at six hours and $-1.6$ below it at twelve on all samples
(above on the deduplicated samples at both, $+3.3$ and $+1.7$), and the time-first encoder
raises the six-hour gain to $+3.3$ on every subset; the hierarchical memory is worth most at
twelve hours ($+3.4$), the horizon at which the morning matters most, and the three seeds of
the twelve-hour model are below the rule on the full stream and above it on independent
samples.  The no-state control is six and a half points above the rule at both horizons
(\cref{tab:exp:nostate:horizons}) and the price history alone already beats the stateful
model ($+1.7$ at 6h, $+5.3$ at 12h): the long-horizon signal is a slow reversal of the 6h and
12h returns, modulated by the price level and by where the market is in its life (time to the
end, age, duration, about 9\% of the gain importance together), which trees read directly and
a direct path that is linear in the returns cannot bend to the level.  The stateful model
with the hand features and the residual construction were not run at these horizons; the
stateless control is the reference there.

\subsection{Synthesis}
\label{sec:horizons:synthesis}

\Cref{tab:horizons:synthesis} and \cref{fig:horizons:synthesis} put the best model of each
family at each horizon on one page.  Three regimes appear.  At one minute the task is activity, calibration is worth five
points to every model, and the sequence of book events carries information about the next
minute of a quiet market that the summary features lose: the stateful model beats calibrated
trees by $+0.6$ alone and $+1.4$ as their residual, with intervals that exclude zero on every
subset, and the gain over the same learner without state is concentrated on the sparse
markets.  At fifteen minutes and one hour the summary is complete: trees, a tabular network
and the stateful model with the same inputs tie, the online regime costs the stateful model
one to two points unless it is given the trees as a base, and the graph and the external
sources add nothing.  At six and twelve hours the task is the direction of a slow reversal
read together with the price level and the market's lifecycle, the stateless control is six
points above the rule and the stateful models of this study are below it.  The same
architecture cannot be the right one at both ends, and the reading the data give is that a
model of this market should be a stateless learner on the activity clock, the first level,
the long returns, the price level and the contract type, with a sequential state added only
where the horizon is shorter than the summary.

\begin{table}[t]
\centering
\caption{The best model of each family at each horizon, all samples and deduplicated
samples (balanced accuracy).  At 1 min, 15 min and 1 h every model carries its own two
cross-fitted logit offsets; at 6 h and 12 h the raw argmax (the stateless control was not
calibrated there and the stateful models' probabilities were not saved).  ``Stateless'' is
LightGBM G5; ``stateful'' is the best run without trees; ``residual'' is the two-pass
stateful model on the cross-fitted trees.  $\Delta$ in points against the stateless control.}
\label{tab:horizons:synthesis}
\small
\begin{fittable}\begin{tabular}{llcccccc}
\toprule
Horizon & Subset & MR & Stateless & Stateful & $\Delta$ & Residual & $\Delta$ \\
\midrule
1 min  & All   & 0.5363 & 0.6623 & 0.6685 & $+0.62$ [$+0.46$, $+0.82$] & \textbf{0.6760} & $+1.37$ [$+1.16$, $+1.58$] \\
       & Dedup & 0.5012 & 0.6143 & 0.6317 & $+1.74$ [$+1.31$, $+2.22$] & \textbf{0.6346} & $+2.04$ [$+1.76$, $+2.30$] \\
15 min & All   & 0.5373 & \textbf{0.6213} & 0.6206 & $-0.07$ [$-0.34$, $+0.19$] & --- & \\
       & Dedup & 0.5253 & \textbf{0.6197} & 0.6158 & $-0.39$ [$-0.63$, $-0.13$] & --- & \\
1 h    & All   & 0.5311 & 0.6044 & 0.5902 & $-1.42$ [$-1.83$, $-1.00$] & \textbf{0.6048} & $+0.04$ [$-0.18$, $+0.24$] \\
       & Dedup & 0.5180 & \textbf{0.5943} & 0.5797 & $-1.46$ [$-1.76$, $-1.14$] & 0.5933 & $-0.10$ [$-0.26$, $+0.05$] \\
6 h    & All   & 0.5069 & \textbf{0.5716} & 0.5396 & $-3.20$ & --- & \\
       & Dedup & 0.5002 & \textbf{0.5665} & 0.5370 & $-2.95$ & --- & \\
12 h   & All   & 0.5071 & \textbf{0.5715} & 0.4914 & $-8.01$ & --- & \\
       & Dedup & 0.4995 & \textbf{0.5680} & 0.5168 & $-5.12$ & --- & \\
\bottomrule
\end{tabular}\end{fittable}
\end{table}

\begin{figure}[t]
\centering
\begin{tikzpicture}
\begin{axis}[thesis, ylabel={$\Delta$ vs the stateless control, points}, symbolic x coords={1 min, 15 min, 1 h, 6 h, 12 h}, xtick=data,
  enlarge x limits=0.14, ymin=-9.8, ymax=3, legend pos=south west, ytick={-8,-6,-4,-2,0,2},
  point meta=explicit symbolic]
\addplot[vizGrid, line width=0.8pt, sharp plot, forget plot] coordinates {(1 min,0) (12 h,0)};
\addplot[only marks, mark=*, mark size=2.4pt, vizA, nodes near coords, every node near coord/.append style={font=\scriptsize, color=black, anchor=north west, xshift=3pt, yshift=-1pt}, error bars/.cd, y dir=both, y explicit, error bar style={vizA, line width=0.9pt}, error mark options={vizA, rotate=90, mark size=2pt}]
  coordinates {(1 min,0.62) +- (0.16,0.20) [$+0.62$] (15 min,-0.07) +- (0.27,0.26) [$-0.07$] (1 h,-1.42) +- (0.41,0.42) [$-1.42$] (6 h,-3.20) [$-3.2$] (12 h,-8.01) [$-8.0$]};
\addplot[only marks, mark=square*, mark size=2.2pt, vizB, nodes near coords, every node near coord/.append style={font=\scriptsize, color=black, anchor=south west, xshift=3pt, yshift=1pt}, error bars/.cd, y dir=both, y explicit, error bar style={vizB, line width=0.9pt}, error mark options={vizB, rotate=90, mark size=2pt}]
  coordinates {(1 min,1.37) +- (0.21,0.21) [$+1.37$] (1 h,0.04) +- (0.22,0.20) [$+0.04$]};
\legend{stateful model, residual of the trees}
\end{axis}
\end{tikzpicture}

\caption{The gain of the stateful family over the stateless control at each horizon
(all samples; calibrated against calibrated at 1 min, 15 min and 1 h with 95\% hour-block
bootstrap intervals; raw argmax at 6 h and 12 h).  The sign reverses with the horizon.}
\label{fig:horizons:synthesis}
\end{figure}

\subsection{Where the models win, by kind of market}
\label{sec:horizons:where}
\label{sec:analysis:where}


The per-market analysis at 1 h (\cref{tab:horizons:where}; the lookback model with the
graph against Mean Reversion on the same 9.86 million samples, before the direct path and the
hand features) shows where a learned model of the book is strong or weak relative to the
rule.  The model is above the rule only on weather markets ($+1.8$) and on markets that last
one hour or less ($+13.7$, a small group of 402 markets in which the rule is near chance);
it is closest to the rule on markets one to six hours from their end date ($+2.4$) and on
the most traded quartile ($-0.3$), and farthest on the fifteen-minute crypto markets
($-6.5$), on sports markets while the game is being played ($-6.6$), on political markets
($-4.3$) and on markets that had no trades in the training days ($-5.3$).  By price level
(\cref{fig:horizons:pricelevel}) both the model and the rule are best between $0.05$ and
$0.2$ and between $0.8$ and $0.95$, and the model collapses above $0.95$ ($-10.6$), where the
process is a one-tick oscillation that the rule reads directly.  The pattern has one reading:
the model wins where the book moves for reasons the book shows (weather resolves on a
schedule, a short market lives inside its own activity), and loses where the price is driven
by an external reference the model does not see, the exchange price for the crypto markets
\cite{young2026openmarket}, the game for in-play sports, the news for politics.  Every
later model of the thesis inherits this map; the hand features and the trees raise every
row, and the control's gain over the rule is positive in every category
(\cref{sec:stateless:control}), but the ordering of the categories, crypto and in-play
sports hardest, does not change.

\begin{table}[t]
\centering
\caption{Balanced accuracy at 1 h by kind of market: the lookback stateful model (graph
variant, seed 0) against Mean Reversion on the same validation samples, with the class mix
of each group.  $\Delta$ in points.}
\label{tab:horizons:where}
\small
\begin{fittable}\begin{tabular}{lrrccc}
\toprule
Group & samples & markets & \textsc{const} share & Model / MR & $\Delta$ \\
\midrule
\multicolumn{6}{l}{\emph{By category}} \\
politics          & 3{,}762{,}863 & 4{,}502 & 0.64 & 0.500 / 0.543 & $-4.3$ \\
sports            & 3{,}088{,}176 & 6{,}720 & 0.54 & 0.512 / 0.535 & $-2.3$ \\
culture           & 1{,}404{,}248 &   726   & 0.37 & 0.474 / 0.490 & $-1.6$ \\
economy, tech     &   889{,}481   & 1{,}197 & 0.55 & 0.470 / 0.498 & $-2.8$ \\
crypto            &   292{,}307   & 1{,}160 & 0.56 & 0.461 / 0.526 & $-6.5$ \\
weather           &   200{,}158   &   696   & 0.30 & 0.480 / 0.463 & $+1.8$ \\
esports           &   150{,}765   & 1{,}213 & 0.33 & 0.504 / 0.517 & $-1.3$ \\
\midrule
\multicolumn{6}{l}{\emph{By time to the end date}} \\
less than 1 h     &    49{,}044   &   539   & 0.09 & 0.371 / 0.425 & $-5.4$ \\
1 h to 6 h        &   136{,}031   & 1{,}250 & 0.25 & 0.490 / 0.466 & $+2.4$ \\
6 h to 1 day      &   219{,}307   & 1{,}615 & 0.55 & 0.472 / 0.471 & $+0.1$ \\
1 day to 1 week   & 2{,}260{,}895 & 2{,}798 & 0.44 & 0.483 / 0.505 & $-2.2$ \\
more than 1 month & 4{,}914{,}916 & 7{,}711 & 0.65 & 0.496 / 0.532 & $-3.7$ \\
\midrule
\multicolumn{6}{l}{\emph{By market duration}} \\
1 h or less       &    24{,}215   &   402   & 0.15 & 0.514 / 0.376 & $+13.7$ \\
1 h to 1 day      &   413{,}522   & 1{,}466 & 0.17 & 0.466 / 0.479 & $-1.3$ \\
more than 3 months & 5{,}954{,}469 & 8{,}484 & 0.61 & 0.501 / 0.531 & $-3.0$ \\
\midrule
\multicolumn{6}{l}{\emph{Sports, by phase}} \\
before the game   & 1{,}160{,}800 & 3{,}932 & 0.63 & 0.467 / 0.494 & $-2.8$ \\
game started      &   557{,}082   & 1{,}818 & 0.11 & 0.464 / 0.530 & $-6.6$ \\
\midrule
\multicolumn{6}{l}{\emph{By traded volume in the training days}} \\
top quartile      & 3{,}071{,}194 & 2{,}291 & 0.48 & 0.510 / 0.512 & $-0.3$ \\
bottom quartile   & 1{,}373{,}930 & 2{,}159 & 0.65 & 0.493 / 0.543 & $-4.9$ \\
no trades         & 1{,}987{,}456 & 7{,}179 & 0.49 & 0.499 / 0.552 & $-5.3$ \\
\bottomrule
\end{tabular}\end{fittable}
\end{table}

\begin{figure}[t]
\centering
\begin{tikzpicture}
\begin{axis}[thesis, ylabel={balanced accuracy at 1 h}, symbolic x coords={$<0.05$, $0.05$--$0.2$, $0.2$--$0.4$, $0.4$--$0.6$, $0.6$--$0.8$, $0.8$--$0.95$, $>0.95$},
  xtick=data, xlabel={mid price at the sample time}, ymin=0.40, ymax=0.65, legend pos=north west, enlarge x limits=0.1]
\addplot[vizA, mark=*, mark size=1.8pt] coordinates {($<0.05$,0.477) ($0.05$--$0.2$,0.534) ($0.2$--$0.4$,0.494) ($0.4$--$0.6$,0.468) ($0.6$--$0.8$,0.476) ($0.8$--$0.95$,0.557) ($>0.95$,0.430)};
\addplot[vizB, mark=square*, mark size=1.8pt] coordinates {($<0.05$,0.521) ($0.05$--$0.2$,0.552) ($0.2$--$0.4$,0.525) ($0.4$--$0.6$,0.508) ($0.6$--$0.8$,0.503) ($0.8$--$0.95$,0.572) ($>0.95$,0.535)};
\legend{stateful model (lookback), Mean Reversion}
\end{axis}
\end{tikzpicture}

\caption{Balanced accuracy at 1 h by price level: the lookback stateful model against
Mean Reversion on the same samples.  Both are best just inside the boundaries and the model
collapses above $0.95$.}
\label{fig:horizons:pricelevel}
\end{figure}

\subsection{Robustness}
\label{sec:horizons:robustness}

\paragraph{Seeds.}  Three seeds of the best LOB-only configuration give $0.5547$, $0.5546$
and $0.5549$ at one hour, a spread of three ten-thousandths, and the time-first encoder and
the forecast features reproduce their gains on a second seed to within a tenth of a point
(\cref{sec:stateful:chain,sec:stateful:encoders,sec:external:tsfm}); the tabular network
is reported as the mean of three seeds.  At twelve hours the seeds spread by $0.009$ and
straddle the rule on the full stream, which is why the twelve-hour result is stated as
``real on independent samples, fragile on the full stream''.  The hour-block bootstrap
intervals of every difference are the second guard: a difference whose interval includes
zero is reported as a tie.

\paragraph{Calibration.}  The two cross-fitted logit offsets of \cref{sec:data:protocol}
are worth $+0.3$ points to the trees at one hour and $+5.4$ at one minute, and the prior
correction of \textcite{menon2021logitadjustment} is destructive on a focal-loss model
(\cref{sec:stateful:hand}); every ordering reported in this section survives the
calibration, several gaps do not, and the one-hour tie between trees and the residual
stateful model only appears once both are calibrated (\cref{sec:stateful:calibration}).

\paragraph{Two views of every sample.}  Every number is given on the full stream and on
the deduplicated subset, and the two disagree in a systematic way: the full stream rewards
what is right on the busy markets, the deduplicated subset what is right on the sparse ones.
The time-first encoder and the forecast features gain on the first; the hand features, the
second pass and the residual construction on the second.  A result that holds on one view
only is reported as such.

\subsubsection{Short tests versus full-length runs}
\label{sec:horizons:short-tests}

A recurring methodological lesson: a two-day short test is not a reliable predictor of
a seven-day result.  The Mean-Reversion initialisation of the direct path gained two
points in the short test and lost two at full length; the five-lookback and global-state
additions kept their sign.  The explanation is that a short test is in the regime where
the recurrent part has learned little and the heuristic is all there is; the full run
is in the regime where the residual has been learned.  Every claim of the thesis rests on
full-length runs, and the short tests are reported as what they are: screening.
